主比较从主机内存中的FP32查询开始，到流完成后top-10 ID与分数回到主机为止
（\cref{fig:3-boundary}）。区间包含pinned staging、H2D、查询预处理、搜索和D2H。
CAGRA int8查询量化计入其中，量化器训练和索引构建在区间外。
资源句柄与输出缓冲复用。JHQ回传到pageable数组，基线回传到pinned数组；
CAGRA/PQ经Python API，JHQ/RaBitQ经C++。解释微小差距时保留这些接口差异。

\par\medskip\noindent\begin{minipage}{\linewidth}
\centering\captionsetup{hypcap=false}
\resizebox{\linewidth}{!}{\begin{tikzpicture}[
 font=\small\linespread{1}\selectfont,
 box/.style={draw=ThesisNavy!65,rounded corners=2pt,align=center,
 text width=2.55cm,minimum height=1.55cm,inner sep=4pt},
 arr/.style={-{Latex[length=2mm]},draw=ThesisNavy!80,thick}]
\node[box] (host) at (0,0) {主机FP32查询\\Pinned staging};
\node[box,fill=ThesisNavy!7] (upload) at (3.5,0) {H2D传输\\查询预处理};
\node[box,fill=ThesisNavy!7] (gpu) at (7,0) {GPU搜索\\选择／精排};
\node[box] (result) at (10.5,0) {D2H ID与分数\\等待流完成};
\draw[arr] (host) -- (upload);
\draw[arr] (upload) -- (gpu);
\draw[arr] (gpu) -- (result);
\draw[ThesisNavy,thick] (-1.4,1.15) -- (-1.4,1.45) -- (11.9,1.45) -- (11.9,1.15);
\node[fill=white,inner sep=3pt] at (5.25,1.45) {主口径：主机到主机计时区间};
\node[box,text width=5.7cm,minimum height=1.3cm] at (1.65,-2.35)
 {搜索计时之外\\加载数据／训练或加载模型／add／预热};
\node[box,text width=5.7cm,minimum height=1.3cm] at (8.65,-2.35)
 {独立策略成本\\预算校准；基于模型的复用回本};
\draw[densely dashed,ThesisNavy!60] (-1.4,-1.1) -- (11.9,-1.1);
\end{tikzpicture}
}
\captionof{figure}[主机与设备计时边界]{计时边界。主机区间包含传输和查询相关计算；设备阶段profile解释其组成。
单独profile与同步可能改变执行，因此阶段时间之和不能替代端到端延迟。}
\label{fig:3-boundary}
\end{minipage}\par\medskip

第$r$轮有$m_r$次计时调用，每次回答$b$条查询：
\begin{equation}
Q_r=\frac{m_rb}{\sum_{i=1}^{m_r}t_{ri}},\qquad
Q=\operatorname{median}_r Q_r,
\label{eq:3-throughput}
\end{equation}
其中$t$以秒计。E1采用$b=1{,}000$和5轮测量。
JHQ每个变体每轮先做1次不计时调用，再通常计时24次；基线先做3次初始预热，
每轮通常计时24次，对判定较慢的点减少次数，并保留实际次数。
预热使JIT和graph capture不进入稳态计时。batch调用的中位／p95延迟描述重复
batch调用，不是逐查询或并发请求的延迟分布。

各轮吞吐的中位数与单次调用延迟中位数的倒数不同。在可得时同时保留二者、
各轮最小／最大值及基线逐调用轨迹。有效慢调用保留在主估计中。
各轮共享训练索引，预算抽样共享查询池：其次数衡量计时波动或抽样敏感性，
不等于独立训练重复数。微小配对差异结合离散度和一致性解释，
不对扫描单元附加依赖独立性假设的显著性主张。

主环境为RTX 5090，可报告显存32,607\,MiB，驱动595.71.05，
CUDA toolkit 13.0和cuVS 26.08.01。运行获取协作GPU锁并保存环境快照；
25核容器配额内使用16个OpenMP线程。E1记录功率上限575\,W，
快照并不证明锁频或恒定宿主负载。JHQ Release的compute\_80 PTX由驱动JIT\cite{cuda_platform}，
库基线使用安装的构建。辅助设备常驻计时与阶段profile分别标注，
因为JHQ主API接受主机查询。

显存核算区分持久编码数组、模型／IVF元数据、查询工作区和运行时／context。
进程used样本包含索引以外的分配；构建前后差值不包含之后创建的内存池。
E1中减context后的JHQ总量与CAGRA索引差值保留各自定义。
组件模型解释常驻总量并显示未归因余量（\cref{sec:4-4,sec:7-7}），
不能代替等显存性能比较；补证以20\,ms设备采样另测构建峰值。冷train/add计时同样不含输入加载、
此前的对象构造与主机模型设置。

\paragraph{补证的测量边界。}
补证实验新增的检查见\cref{sec:a-supplement}。硬件型号及报告驱动／toolkit与归档一致，
但计时属于另一会话。算术模式及缓存模型下的精度干预见\cref{sec:4-2,sec:7-precision}；
历史基线比较保留会话区别。
